\subsection{Extraction of Blink Regions}
In Stage~C, the threshold is scanned across the signal in
the manner of the Kleifges crossing detector: maximal contiguous runs of supra-threshold
samples define candidate blink regions, brief crossings shorter than a minimum-duration gate
are discarded, and the accepted regions are mapped from their concatenated coordinates back
to epoch-relative onset and offset times.

The final stage applies the sample-level thresholds
\(\Theta=\{\theta_c\}_{c=1}^{C}\) to obtain channel-wise blink regions. For
each channel \(c\), the epochs selected for localization are ordered as
\(\mathcal{S}=\{i_1,\ldots,i_R\}\), where \(R=|\mathcal{S}|\). Their samples
are concatenated into a one-dimensional signal
\[
u_c[n] = x_{i_r,c,t},
\qquad
n=(r-1)T+t,\quad
r \in \{1,\ldots,R\},\quad
t \in \{1,\ldots,T\}.
\]
Thus \(u_c \in \mathbb{R}^{RT}\) is the channel-specific signal on which blink
onsets and offsets are estimated.

%If the analysis is run in the fallback mode  described above, \(\mathcal{S}\) is the full valid epoch set; otherwise it is the set of epochs identified as suspicious in Stage A.

Candidate blink samples are defined by a threshold-crossing indicator
\[
h_c[n] = \mathbf{1}\{u_c[n] > \theta_c\},
\qquad n=1,\ldots,RT .
\]
After any required polarity normalization, \(h_c[n]=1\) indicates that the
sample amplitude exceeds the channel-specific blink-region threshold. Blink
regions are then defined as maximal contiguous runs of supra-threshold samples.
A candidate interval for channel \(c\) is written as the half-open index range
\([s_{c,r},e_{c,r})\), where
\[
h_c[n]=1 \quad \text{for all } n=s_{c,r},\ldots,e_{c,r}-1,
\]
with either \(s_{c,r}=1\) or \(h_c[s_{c,r}-1]=0\), and either
\(e_{c,r}=RT+1\) or \(h_c[e_{c,r}]=0\). Equivalently, the onset
\(s_{c,r}\) is the first sample in a run that crosses above \(\theta_c\), and
the offset \(e_{c,r}\) is the first subsequent sample at which the signal has
returned to the sub-threshold state.

%To suppress brief threshold crossings that are unlikely to represent complete
%blinks, each candidate interval is required to exceed a minimum duration. Let
%\(f_s\) be the sampling frequency and let \(\ell_{\min}\) be the minimum blink
%duration in seconds. The corresponding sample length is
%\[
%L_{\min}=\ell_{\min} f_s .
%\]
%The accepted blink-region set for channel \(c\) is
%\[
%\mathcal{R}_c
%=
%\left\{
%(s_{c,r},e_{c,r}) :
%e_{c,r}-s_{c,r} > L_{\min}
%\right\}.
%\]
%This duration gate is applied independently for each channel. Consecutive
%supra-threshold samples are merged into the same region by construction, while
%two runs separated by at least one sub-threshold sample are treated as distinct
%candidate blink regions.

Finally, the accepted intervals are mapped from concatenated coordinates back to epoch-local coordinates using the known epoch boundaries.

%If an accepted
%onset \(s_{c,r}\) lies in the concatenated segment corresponding to epoch
%\(i_q\), its within-epoch onset sample is
%\[
%\hat{s}_{c,r}=s_{c,r}-(q-1)T,
%\]
%and its onset time is \((\hat{s}_{c,r}-1)/f_s\) seconds relative to the
%beginning of that epoch. The interval duration is
%\[
%\hat{d}_{c,r}=\frac{e_{c,r}-s_{c,r}}{f_s}.
%\]
The final output of this stage is therefore a channel-indexed collection of
blink regions, each represented by an epoch identifier, a channel label, an
onset index or time, and an offset or duration.

%The extraction remains
%channel-wise throughout: the threshold \(\theta_c\) estimated for channel \(c\)
%is applied only to that channel's signal, and the resulting regions
%\(\mathcal{R}_c\) are retained as channel-specific blink candidates.
